\paragraph{全息 Bernoulli 测度的不动点、自相似质量与精确维数}
在全息地址像的自仿射动力学、圆柱--球严格同构以及 Bernoulli 推前测度的 exact-dimensional 性都已建立之后，还可把这一地址系统进一步压缩为以下三条可直接调用的测度闭式。

\begin{proposition}[\texorpdfstring{$2^L$}{2^L}-进全息点的仿射 IFS 固定点方程]\label{prop:xi-time-part62dca-hologram-affine-ifs-fixedpoint-equation}
\leanverified{paper\_xi\_time\_part62dca\_hologram\_affine\_ifs\_fixedpoint\_equation}
沿用定理 \ref{thm:xi-time-part62dc-hologram-selfaffine-cantor-dynamics} 的记号。设 \(E\) 为有限事件集，
\[
\mathrm{code}:E\hookrightarrow \NN,
\qquad
B=2^L>\max_{a\in E}\mathrm{code}(a),
\qquad
T:=T_2(E_{\min}),
\]
并记
\[
X(\mathbf e)=\sum_{t\ge 1}\mathrm{code}(e_t)B^{t-1}v
\qquad
\bigl(\mathbf e=(e_t)_{t\ge 1}\in E^{\NN}\bigr).
\]
对每个 \(a\in E\)，定义仿射映射
\[
F_a:T\longrightarrow T,
\qquad
F_a(x):=\mathrm{code}(a)\,v+B x.
\]
再取任意严格正概率向量
\[
\pi=(\pi_a)_{a\in E},
\]
并记 \(\mathbf P_\pi\) 为 \(E^{\NN}\) 上的 Bernoulli 乘积测度，
\[
\mu_\pi:=X_*\mathbf P_\pi.
\]
则成立：
\begin{enumerate}
  \item \(\mu_\pi\) 满足精确不动点方程
        \[
        \mu_\pi=\sum_{a\in E}\pi_a\,(F_a)_*\mu_\pi;
        \]
  \item 对任意有限词
        \[
        \mathbf a=(a_1,\dots,a_n)\in E^n,
        \]
        若记
        \[
        U_{\mathbf a}:=F_{a_1}\circ\cdots\circ F_{a_n}(T),
        \]
        则
        \[
        \mu_\pi(U_{\mathbf a})=\pi_{a_1}\cdots \pi_{a_n}.
        \]
\end{enumerate}
\end{proposition}
\begin{proof}
对任意 \(a\in E\) 与任意尾流 \(\mathbf e'=(e_2,e_3,\dots)\in E^{\NN}\)，由地址级数的首项分裂式直接得到
\[
X(a,\mathbf e')
=
\mathrm{code}(a)\,v+B X(\mathbf e')
=
F_a(X(\mathbf e')).
\]
按第一位事件对 \(\mathbf P_\pi\) 作条件分解，便有
\[
\mu_\pi
=
X_*\mathbf P_\pi
=
\sum_{a\in E}\pi_a\,(F_a)_*X_*\mathbf P_\pi
=
\sum_{a\in E}\pi_a\,(F_a)_*\mu_\pi,
\]
从而第 \(1\) 条成立。

再对 \(\mathbf a=(a_1,\dots,a_n)\in E^n\)，记 cylinder 集
\[
[\mathbf a]
:=
\{\mathbf e\in E^{\NN}:\ e_1=a_1,\dots,e_n=a_n\}.
\]
由上式迭代可得
\[
X([\mathbf a])=U_{\mathbf a}\cap X(E^{\NN}).
\]
另一方面，\(\mathbf P_\pi([\mathbf a])=\pi_{a_1}\cdots \pi_{a_n}\)。因此
\[
\mu_\pi(U_{\mathbf a})
=
\mu_\pi\!\bigl(U_{\mathbf a}\cap X(E^{\NN})\bigr)
=
\mathbf P_\pi([\mathbf a])
=
\pi_{a_1}\cdots \pi_{a_n}.
\]
第 \(2\) 条成立。证毕。
\end{proof}

\leanverified{paper\_xi\_time\_part62dca\_hologram\_bernoulli\_exact\_dimension\_formula}
\begin{theorem}[全息 Bernoulli 测度的精确维数公式]\label{thm:xi-time-part62dca-hologram-bernoulli-exact-dimension-formula}
沿用命题 \ref{prop:xi-time-part62dca-hologram-affine-ifs-fixedpoint-equation} 的记号，并记
\[
H(\pi):=-\sum_{a\in E}\pi_a\log \pi_a.
\]
则 \(\mu_\pi\) 为 exact-dimensional 测度，并且对 \(\mu_\pi\)-几乎处处的
\[
x\in X(E^{\NN})
\]
都有
\[
\lim_{r\downarrow 0}\frac{\log \mu_\pi(B_T(x,r))}{\log r}
=
\frac{H(\pi)}{\log B},
\]
其中 \(B_T(x,r)\) 表示 \(T\) 上的 \(B\)-进超度量球。特别地，
\[
\dim_{\mathrm H}(\mu_\pi)=\frac{H(\pi)}{\log B}.
\]
\end{theorem}
\begin{proof}
取 \(x=X(\mathbf e)\)，其中
\[
\mathbf e=(e_1,e_2,\dots)\in E^{\NN}.
\]
对每个 \(n\ge 1\)，记前缀
\[
\mathbf a_n:=(e_1,\dots,e_n).
\]
由结论 \ref{con:xi-time-part62dgb-hologram-cylinder-ball-isometry}，半径 \(B^{-n}\) 的球恰对应于长度 \(n\) 的 cylinder 像：
\[
B_T(x,B^{-n})
=
U_{\mathbf a_n},
\qquad
U_{\mathbf a_n}\cap X(E^{\NN})=X([\mathbf a_n]).
\]
于是由命题 \ref{prop:xi-time-part62dca-hologram-affine-ifs-fixedpoint-equation}，
\[
\mu_\pi\!\bigl(B_T(x,B^{-n})\bigr)
=
\mu_\pi(U_{\mathbf a_n})
=
\prod_{t=1}^{n}\pi_{e_t}.
\]
从而
\[
\frac{\log \mu_\pi(B_T(x,B^{-n}))}{\log B^{-n}}
=
\frac{-\frac1n\sum_{t=1}^{n}\log \pi_{e_t}}{\log B}.
\]
在 \(\mathbf P_\pi\) 下，\((e_t)\) 为独立同分布序列，因此由强大数定律，
\[
\frac1n\sum_{t=1}^{n}(-\log \pi_{e_t})\longrightarrow H(\pi)
\qquad \mathbf P_\pi\text{-几乎处处}.
\]
经 \(X\) 推前后即得
\[
\lim_{n\to\infty}
\frac{\log \mu_\pi(B_T(x,B^{-n}))}{\log B^{-n}}
=
\frac{H(\pi)}{\log B}
\qquad \mu_\pi\text{-几乎处处}.
\]
又由于 \(T\) 上的该超度量具有离散尺度结构，对每个充分小的半径 \(r\) 都存在唯一的 \(n\) 使相应球与某个半径 \(B^{-n}\) 的球一致，故上式等价于
\[
\lim_{r\downarrow 0}\frac{\log \mu_\pi(B_T(x,r))}{\log r}
=
\frac{H(\pi)}{\log B}.
\]
因此 \(\mu_\pi\) exact-dimensional，且其 Hausdorff 维数即为该常数。证毕。
\end{proof}

\leanverified{paper\_xi\_time\_part62dca\_hologram\_bernoulli\_lq\_spectrum}
\begin{theorem}[全息 Bernoulli 推前测度的精确 \texorpdfstring{$L^q$}{Lq}-谱]\label{thm:xi-time-part62dca-hologram-bernoulli-lq-spectrum}
沿用定理 \ref{thm:xi-time-part62dca-hologram-bernoulli-exact-dimension-formula} 的记号，并额外假设
\[
\pi_a>0
\qquad
(a\in E).
\]
对任意 \(q\in\RR\)，定义第 \(n\) 层 cylinder 质量和
\[
Z_n(q):=\sum_{\mathbf a\in E^n}\mu_\pi(U_{\mathbf a})^q.
\]
则极限
\[
\tau_{\mu_\pi}(q)
:=
\lim_{n\to\infty}\frac{\log Z_n(q)}{\log B^{-n}}
\]
存在，并满足显式公式
\[
\boxed{\;
\tau_{\mu_\pi}(q)
=
-\frac{\log\sum_{a\in E}\pi_a^q}{\log B}
=
-\frac{\log\sum_{a\in E}\pi_a^q}{L\log 2}.\;}
\]
特别地，\(\mu_\pi\) 的 exact-dimensional 性与维数公式
\[
\dim_{\mathrm H}(\mu_\pi)=\frac{H(\pi)}{\log B}
\]
正是这一地址模型在 \(q=1\) 处的熵极限读出。
\end{theorem}
\begin{proof}
由命题 \ref{prop:xi-time-part62dca-hologram-affine-ifs-fixedpoint-equation}，对任意
\[
\mathbf a=(a_1,\dots,a_n)\in E^n
\]
都有
\[
\mu_\pi(U_{\mathbf a})=\pi_{a_1}\cdots \pi_{a_n}.
\]
因此
\[
Z_n(q)
=
\sum_{(a_1,\dots,a_n)\in E^n}(\pi_{a_1}\cdots \pi_{a_n})^q
=
\left(\sum_{a\in E}\pi_a^q\right)^n.
\]
于是
\[
\frac{\log Z_n(q)}{\log B^{-n}}
=
\frac{n\log\sum_{a\in E}\pi_a^q}{-n\log B}
=
-\frac{\log\sum_{a\in E}\pi_a^q}{\log B}.
\]
右端与 \(n\) 无关，故极限存在并等于所示常数。再由
\[
\log B=L\log 2
\]
得到第二种写法。

最后，定理 \ref{thm:xi-time-part62dca-hologram-bernoulli-exact-dimension-formula} 已经给出
\[
\dim_{\mathrm H}(\mu_\pi)=\frac{H(\pi)}{\log B},
\]
故这里的 \(\tau\)-公式把此前的 exact-dimensional 结论与全体 \(q\)-矩谱统一到同一个闭式中。证毕。
\end{proof}

\begin{corollary}[均匀全息像的 Ahlfors 正则性与精确 Hausdorff 量级]\label{cor:xi-time-part62dca-hologram-uniform-ahlfors-regularity}
\leanverified{paper\_xi\_time\_part62dca\_hologram\_uniform\_ahlfors\_regularity}
在定理 \ref{thm:xi-time-part62dca-hologram-bernoulli-exact-dimension-formula} 的设定下，取均匀分布
\[
\pi_a\equiv |E|^{-1},
\qquad
s:=\frac{\log |E|}{\log B},
\qquad
\mathcal X:=X(E^{\NN})\subset T.
\]
则成立：
\begin{enumerate}
  \item 对任意 \(\mathbf a\in E^n\)，有
        \[
        \mu_{\mathrm{unif}}(U_{\mathbf a})
        =
        |E|^{-n}
        =
        B^{-ns};
        \]
  \item 对任意 \(x\in\mathcal X\) 与任意 \(n\ge 0\)，有
        \[
        \mu_{\mathrm{unif}}(B_T(x,B^{-n}))
        =
        B^{-ns};
        \]
  \item \(\mathcal X\) 是 \(s\)-Ahlfors 正则的，并且
        \[
        0<\mathcal H^s(\mathcal X)<\infty.
        \]
\end{enumerate}
\end{corollary}
\begin{proof}
第 \(1\) 条由命题 \ref{prop:xi-time-part62dca-hologram-affine-ifs-fixedpoint-equation} 直接得到，因为均匀分布下
\[
\pi_{a_1}\cdots \pi_{a_n}=|E|^{-n}=B^{-ns}.
\]

再由结论 \ref{con:xi-time-part62dgb-hologram-cylinder-ball-isometry}，\(\mathcal X\) 上半径 \(B^{-n}\) 的球恰等于某个 \(U_{\mathbf a}\)，其与 \(\mathcal X\) 的交正对应于长度 \(n\) 的 cylinder 像。故第 \(2\) 条立即由第 \(1\) 条推出。

第 \(2\) 条表明，在全部离散尺度 \(r=B^{-n}\) 上都存在精确体积律
\[
\mu_{\mathrm{unif}}(B_T(x,r))=r^s
\qquad (x\in\mathcal X).
\]
这正是超度量口径下的 \(s\)-Ahlfors 正则性。于是由标准质量分布原理与覆盖估计，得到
\[
0<\mathcal H^s(\mathcal X)<\infty.
\]
证毕。
\end{proof}

\begin{corollary}[均匀全息像上的 Hausdorff 测度识别]\label{cor:xi-time-part62dca-hologram-hausdorff-measure-identification}
沿用推论 \ref{cor:xi-time-part62dca-hologram-uniform-ahlfors-regularity} 的记号，记
\[
\mathcal X:=X(E^{\NN}),
\qquad
\nu_s:=\mathcal H^s\!\restriction_{\mathcal X}.
\]
则
\[
\mathcal H^s(\mathcal X)=1,
\qquad
\nu_s=\mu_{\mathrm{unif}}.
\]
特别地，
\[
\mu_{\mathrm{unif}}
=
\frac{\mathcal H^s\!\restriction_{\mathcal X}}{\mathcal H^s(\mathcal X)}.
\]
\end{corollary}
\begin{proof}
若 \(|E|=1\)，则 \(\mathcal X\) 为单点集、\(s=0\)，而 \(\mu_{\mathrm{unif}}\) 与 \(\mathcal H^0\!\restriction_{\mathcal X}\) 都是该点上的 Dirac 概率测度，结论显然成立。
以下设 \(|E|\ge 2\)。

由结论 \ref{con:xi-time-part62dgb-hologram-cylinder-ball-isometry}，\(\mathcal X\) 中半径 \(B^{-n}\) 的闭球恰对应于长度 \(n\) 的 cylinder 像 \(U_{\mathbf a}\)。
再由推论 \ref{cor:xi-time-part62dca-hologram-uniform-ahlfors-regularity}，
\[
\mu_{\mathrm{unif}}(U_{\mathbf a})=B^{-ns}.
\]
由于 \(|E|\ge 2\)，每个长度 \(n\) cylinder 都含有两条在第 \(n+1\) 位分歧的延拓，故
\[
\operatorname{diam}(U_{\mathbf a})=B^{-n},
\qquad
\mu_{\mathrm{unif}}(U_{\mathbf a})=\operatorname{diam}(U_{\mathbf a})^s.
\]

固定任意 cylinder \(U_{\mathbf a}\)。
在当前离散超度量口径下，\(\mathcal H^s\) 可由球覆盖计算。
若 \(U_{\mathbf a}\subset \bigcup_i V_i\) 是任意由 \(\mathcal X\) 中闭球组成的可数覆盖，则每个 \(V_i\) 也是某个 cylinder，故
\[
\mu_{\mathrm{unif}}(U_{\mathbf a})
\le
\sum_i \mu_{\mathrm{unif}}(V_i)
=
\sum_i \operatorname{diam}(V_i)^s.
\]
对所有直径受控的球覆盖取下确界，得到
\[
\mathcal H^s(U_{\mathbf a})\ge \mu_{\mathrm{unif}}(U_{\mathbf a}).
\]

反过来，对任意 \(m\ge |\mathbf a|\)，\(U_{\mathbf a}\) 可由其所有长度 \(m\) 的后继 cylinder 两两不交覆盖；这类后继共有 \(|E|^{m-|\mathbf a|}\) 个，每个直径都是 \(B^{-m}\)。于是
\[
\mathcal H^s(U_{\mathbf a})
\le
|E|^{m-|\mathbf a|}B^{-ms}
=
B^{-|\mathbf a|s}
=
\mu_{\mathrm{unif}}(U_{\mathbf a}).
\]
故
\[
\mathcal H^s(U_{\mathbf a})=\mu_{\mathrm{unif}}(U_{\mathbf a})
\]
对每个 cylinder 都成立。

取空前缀对应的 \(U_{\emptyset}=\mathcal X\)，便得
\[
\mathcal H^s(\mathcal X)=\mu_{\mathrm{unif}}(\mathcal X)=1.
\]
最后，cylinder 族构成 \(\mathcal X\) 的可数 clopen \(\pi\)-系统，并生成其 Borel \(\sigma\)-代数。
有限 Borel 测度 \(\nu_s\) 与 \(\mu_{\mathrm{unif}}\) 在全部 cylinder 上一致，故由测度唯一性，
\[
\nu_s=\mu_{\mathrm{unif}}.
\]
归一化公式随即成立。证毕。
\end{proof}

\begin{corollary}[等概率全息地址测度的 Rényi 维数与多重分形谱完全塌缩]\label{cor:xi-time-part62dca-hologram-uniform-monofractal-collapse}
\leanverified{paper\_xi\_time\_part62dca\_hologram\_uniform\_monofractal\_collapse}
沿用推论 \ref{cor:xi-time-part62dca-hologram-hausdorff-measure-identification} 的记号，并设
\[
s:=\frac{\log |E|}{\log B}
=
\frac{\log |E|}{L\log 2}.
\]
则均匀测度 \(\mu_{\mathrm{unif}}\) 满足
\[
\tau_{\mu_{\mathrm{unif}}}(q)=(q-1)s
\qquad
(q\in\RR).
\]
若对 \(q\neq 1\) 定义 Rényi 维数
\[
D_q(\mu_{\mathrm{unif}}):=\frac{\tau_{\mu_{\mathrm{unif}}}(q)}{q-1},
\]
并在 \(q=1\) 处取连续延拓，则
\[
\boxed{\;
D_q(\mu_{\mathrm{unif}})=s
\qquad
(q\in\RR).\;}
\]
进一步，\(\mu_{\mathrm{unif}}\) 的多重分形谱塌缩为单点：
\[
\boxed{\;
f(\alpha)=
\begin{cases}
s,& \alpha=s,\\
-\infty,& \alpha\neq s.
\end{cases}\;}
\]
\end{corollary}
\begin{proof}
把均匀概率
\[
\pi_a=|E|^{-1}
\]
代入定理 \ref{thm:xi-time-part62dca-hologram-bernoulli-lq-spectrum}，得到
\[
\sum_{a\in E}\pi_a^q
=
|E|\cdot |E|^{-q}
=
|E|^{1-q}.
\]
故
\[
\tau_{\mu_{\mathrm{unif}}}(q)
=
-\frac{\log |E|^{1-q}}{\log B}
=
\frac{(q-1)\log|E|}{\log B}
=(q-1)s.
\]
于是对 \(q\neq 1\)，
\[
D_q(\mu_{\mathrm{unif}})
=
\frac{\tau_{\mu_{\mathrm{unif}}}(q)}{q-1}
=
s.
\]
而 \(q\to 1\) 的连续延拓同样给出 \(D_1=s\)。

另一方面，由推论 \ref{cor:xi-time-part62dca-hologram-uniform-ahlfors-regularity}，
\[
\mu_{\mathrm{unif}}(B_T(x,B^{-n}))=B^{-ns}
\qquad
(x\in \mathcal X,\ n\ge 0),
\]
故每一点的局部维数都等于 \(s\)。因此
\[
\{x:\dim_{\mathrm{loc}}(\mu_{\mathrm{unif}},x)=s\}=\mathcal X,
\]
而对 \(\alpha\neq s\)，相应层集为空。再由推论
\ref{cor:xi-time-part62dca-hologram-hausdorff-measure-identification}
中的
\[
\mathcal H^s(\mathcal X)=1,
\]
得到该唯一非空层集的 Hausdorff 维数正是 \(s\)。这就给出所示的单点多重分形谱。证毕。
\end{proof}

\noindent
因此，\(2^L\)-进全息地址并不仅把 Bernoulli 时间流嵌入一个 Cantor 型像集，而且把源分布的 Shannon 熵、逐前缀 cylinder 质量以及均匀源下的 Hausdorff 量级同时压缩为完全显式的自相似闭式。

\endinput
